\section{这是译文}
\label{sec:modexpansion}



这是译文~\ref{cor:mod-reduction} 这是译文~\ref{cor:mod-dim-tradeoff} 这是译文~$\calD$ 这是译文~$\Z^{n}$ 这是译文 $(B,\delta)$-这是译文~$B,\delta\geq 0$ 这是译文~$\vec{x} \gets
\calD$ 这是译文 $B$ 这是译文 $\delta$. 

\begin{theorem}
  \label{thm:mod-dim-switch}  这是译文 $m,n, n',q,q' \geq 1$ 这是译文  $\matG \in \Z^{n' \times n}$ 这是译文 $\Lambda =
  \frac{1}{q'} \matG^{T} \Z^{n'} + \Z^{n}$ 这是译文 $\matB$,  这是译文 $\calD$ 这是译文 $(B,\delta)$-这是译文  $\Z^{n}$.  这是译文 $\alpha, \beta > 0$ 这是译文~$\eps \in (0,1/2)$ 这是译文  \[\beta^{2} \geq
  \alpha^{2} + (4/\pi) \ln(2n(1+1/\eps)) \cdot (\max\set{q^{-1}, \length{\gs{\matB}}} \cdot B)^{2}.\]  这是译文 (这是译文) 这是译文~$\lwe_{n,m,q,\le \alpha}(\calD)$ 这是译文 $\lwe_{n',m,q',\leq \beta}(\matG
  \cdot \calD)$ 这是译文 $\delta + 14 \eps m$.
\end{theorem}

这是译文 $\|\widetilde \matB\|$ 这是译文~\ref{lem:gpv}. 这是译文 $\Z_{q'}^{n'}$, 这是译文~\ref{def:lwe}. 这是译文 $\lwe_{n',q',\leq \beta}$  这是译文~$\lwe_{n',q',\beta}$ (这是译文~\ref{lem:lelwe}).

这是译文 (这是译文 \lwe 这是译文), 这是译文 $n'=n$, $\matG=\matI$ 这是译文 $n$-这是译文 $\matB = \matI/q'$.  这是译文 $q \ge q' \ge \sqrt{2\ln(2n(1+1/\eps))} \cdot
(B/\alpha)$ 这是译文 $\beta=\sqrt{2} \alpha$, 这是译文~$B/\alpha$, 这是译文~$\alpha$ 这是译文.

\begin{corollary}
  \label{cor:mod-reduction}  这是译文 $m,n \geq 1$, $q \geq q' \geq 1$, $(B,\delta)$-这是译文~$\calD$ 这是译文~$\Z^{n}$, $\alpha,\beta>0$ 这是译文~$\eps \in (0,1/2)$ 这是译文  \[
	\beta^{2} \geq \alpha^{2} + (4/\pi) \ln(2n(1+1/\eps))
  \cdot (B/q')^2,
	\]  这是译文  $\lwe_{n,m,q,\le \alpha}(\calD)$ 这是译文  $\lwe_{n,m,q',\le \beta}(\calD)$ 这是译文 $\delta + 14 \eps m$.
\end{corollary}

这是译文 \lwe (这是译文~\ref{lem:normalform}), 这是译文 $\calD = D_{\Z^{n}, \sqrt{2} \alpha q}$, 这是译文 (这是译文~\ref{thm:searchtodecisionmicpei}), 这是译文 (这是译文) $\lwe$ 这是译文 \emph{这是译文} 这是译文 $q$. 这是译文 $\calD = D_{\Z^{n}, r}$ 这是译文 $(C r \sqrt{n \log(n/\delta)}, \delta)$-这是译文 $C>0$, 这是译文 $n$ 这是译文. (这是译文~$(r\sqrt{n},2^{-n})$-这是译文~\cite[Lemma~1.5]{banaszczyk93:_new}, 这是译文 $n$.)

\begin{corollary}
  \label{cor:mod-reduction-2}  这是译文 $\delta \in (0,1/2)$, $m \ge n \geq 1$, $q' \ge 25$. 这是译文 $q \in [q',2q')$ 这是译文 $2$ 这是译文 $q'$  这是译文 $\alpha \ge \sqrt{\ln(2n(1+16/\delta)/\pi)}/q$.  这是译文 (这是译文) 这是译文 $\lwe_{n,m,q, \alpha}$ 这是译文 $\lwe_{n,m',q', \le \beta}$  这是译文 $m'=m-(16 n + 4 \ln\ln q)$ 这是译文  \[
	\beta = C \alpha \sqrt{n} \sqrt{\log(n/\delta)\log(m/\delta)}
	\]  这是译文 $C>0$, 这是译文 $\zeta$ 这是译文 $(\zeta-\delta)/4$.
\end{corollary}


  这是译文  $n=kn'$ 这是译文 $k \ge 1$, 这是译文~$q' = q^{k}$.  这是译文~$\vec{G} =
  \vec{I}_{n'} \otimes \vec{g}$, 这是译文~$\vec{g} = (1,q,q^2, \ldots,
  q^{k-1})^T \in \Z^k$.  这是译文 $\Lambda = q^{-k} \vec{G}^T
  \Z^{n'} + \Z^n$.  这是译文~$\Lambda$ 这是译文 \[
\vec{B}= \vec{I}_{n'} \otimes 
\begin{bmatrix}
q^{-1} & q^{-2} & \cdots & q^{-k} \\
       & q^{-1} & \cdots & q^{1-k} \\
       &        & \ddots & \vdots \\
       &		&		 & q^{-1}
     \end{bmatrix} \in \R^{n \times n};
\] 这是译文~$\vec{B}$ 这是译文~$\Lambda$ 这是译文 $\gs{\matB} = q^{-1} \matI$ 这是译文 $\length{\gs{\matB}}=q^{-1}$.  这是译文~$n \log q = n' \log q'$ 这是译文~$\calD = D_{\Z^{n}, \alpha q}$ (这是译文~\ref{lem:normalform}), 这是译文 $(\alpha q
\sqrt{n}, 2^{-n})$-这是译文~$\sqrt{n}$ 这是译文. 

\begin{corollary}
  \label{cor:mod-dim-tradeoff}  这是译文 $n, m,q \ge 1$, $k \ge 1$ 这是译文~$n$,  $(B,\delta)$-这是译文~$\calD$ 这是译文~$\Z^{n}$,  $\alpha,\beta>0$, 这是译文~$\eps \in (0,1/2)$ 这是译文 \[
\beta^{2} \geq \alpha^{2} + (4/\pi)
  \ln(2n(1+1/\eps)) \cdot (B/q)^{2},
\]  这是译文~$\lwe_{n,m,q,\le \alpha}(\calD)$  这是译文~$\lwe_{n/k,m,q^k,\le \beta}(\vec{G} \cdot \calD)$ 这是译文 $\delta+14 \eps m$,  这是译文~$\vec{G} = \vec{I}_{n/k} \otimes  (1,q,q^2, \ldots, q^{k-1})^T$.
\end{corollary}

\noindent 这是译文~\ref{thm:mod-dim-switch} 这是译文.


\begin{lemma}
  \label{lem:reverse-main}  这是译文~\ref{thm:mod-dim-switch}, 这是译文  \[
	r
  \geq \max\set{q^{-1}, \length{\gs{\matB}}} \cdot
  \sqrt{2\ln(2n(1+1/\epsilon))/\pi}. 
	\] 这是译文 $\T_{q}^{n} \times \T$ 这是译文 $\T_{q'}^{n'} \times \T$, 这是译文:
  \begin{itemize}[itemsep=0pt]
  \item 这是译文 $4 \epsilon$ 这是译文.
  \item 这是译文  $A_{q,\vecs,D_{\alpha}}$ 这是译文 $\vecs \in \Z^{n}$ 这是译文  $\length{\vecs} \leq B$, 这是译文 $10 \epsilon$ 这是译文 $A_{q',\matG\vecs,D_{\alpha'}}$, 这是译文  $(\alpha')^{2} = \alpha^{2} + r^{2}(\length{\vecs}^{2}+B^{2}) \leq
    \alpha^{2} + 2(rB)^{2}$. 
%\dnote{I think we use the condition $\length{\vecs} \leq B$ only for that last inequality. 
%If this is indeed the case, I would be in favour of removing the condition, and commenting at the end of this item that 
%if~$\length{\vecs} \leq B$, then $\alpha^{2} + r^{2}(\length{\vecs}^{2}+B^{2}) \leq
%    \alpha^{2} + 2(rB)^{2}$}\onote{don't we also need it when we apply Lemma~\ref{lem:smoothip} towards the end of the proof? (I didn't check too carefully)}
%\dnote{Why? we don't need $\vec{s}$ small too apply it, do we?}\onote{yes, I think we do; look carefully at the lemma...}
  \end{itemize}
\end{lemma}

\begin{proof}
  这是译文 $\T_{q}^{n}$ 这是译文  $\T_{q'}^{n'}$, 这是译文-$1$ 这是译文 $\T_{q}^{n}$ 这是译文 $\vecs \in \Z^{n}$, 这是译文 $\T_{q'}^{n'}$ 这是译文 $\matG \vecs \in \Z^{n'}$, 这是译文 $\veca \in \T_{q}^{n}$ (这是译文) 这是译文 $\veca' \in
  \T_{q'}^{n'}$, 这是译文  \[ \inner{\veca', \matG \vecs} = \inner{\matG^{T} \veca', \vecs}
  \approx \inner{\veca,\vecs} \bmod 1 \] 这是译文 (这是译文) $\vecs \in
  \Z^{n}$.  这是译文 $\veca'$ 这是译文  $\matG^{T} \veca' \approx \veca \bmod \Z^{n}$, 这是译文 $r$.
  

  % In order to formally describe how the reduction samples $\veca'$, we
  % first define an $n$-dimensional lattice \[ \Lambda = q'^{-1} \cdot
  % \lat(\matG^{T}) + \Z^{n} = q'^{-1} \cdot \lat(\matI_{n'} \otimes
  % \vecg) + \Z^{n}, \] and claim that it has as a basis the
  % upper-triangular matrix \[ \matB =
  % \matI_{n'} \otimes
  % \begin{bmatrix}
  %   (q'_{1})^{-1} & \star & \cdots & \star \\
  %   & (q'_{2})^{-1} & \cdots & \star \\
  %   & & \ddots & \vdots \\
  %   & & & (q'_{k})^{-1}
  % \end{bmatrix} \in \R^{n \times n}, \] where the rightmost column of
  % the above displayed $k$-by-$k$ matrix is $q'^{-1} \vecg$, and the
  % $j$th column is obtained by multiplying the $(j+1)$st column by
  % $q'_{j+1}$ and reducing modulo $1$, for $j=k-1,\ldots, 1$.  Indeed,
  % $\matB$ is a basis of $\Lambda$ because all of its columns are in
  % $\Lambda$ by construction, and because
  % $\det(\matB)=q'^{-n'}=\abs{\Lambda/\Z^{n}}^{-n'}=\det(\Lambda)$.  Also
  % notice that because $\matB$ is upper triangular, its Gram-Schmidt
  % orthogonalization (from left to right) is just the diagonal matrix
  % with $(q'_{i})^{-1}$ in the $i$th diagonal, so $\max_{i}
  % \length{\gs{\vecb_{i}}} \leq q_{\min}^{-1}$ and $r \geq
  % \sqrt{2}\smootheps(\Lambda)$ by~\cite[Theorem
  % 3.1]{DBLP:conf/stoc/GentryPV08}\cnote{Put this in prelims...}.
  
  这是译文 $(\veca, b) \in \T_{q}^{n} \times \T$, 这是译文:
  \begin{itemize}
  \item 这是译文 $\vecf \gets D_{\Lambda-\veca,r}$ 这是译文~\ref{lem:gpv} 这是译文  $\matB$, 这是译文 $\vecv = \veca + \vecf \in
    \Lambda/\Z^{n}$. (这是译文 $\Lambda-\veca$ 这是译文 $\veca =
      \bar{\veca}+\Z^{n}$ 这是译文 $\Z^{n} \subseteq \Lambda$.) 这是译文 $\veca' \in \T_{q'}^{n'}$ 这是译文  $\matG^{T} \veca' = \vecv \bmod \Z^{n}$. 这是译文 $\matG^{T}
      \veca' = \veczero \bmod \Z^{n}$, 这是译文  $\matG^{T} \veca' = \vecv \bmod \Z^{n}$.
  \item 这是译文 $e' \gets D_{rB}$ 这是译文 $b' = b + e' \in \T$.  
  \item 这是译文 $(\veca',b')$.
  \end{itemize}

  这是译文 $\veca'$ 这是译文  (这是译文) 这是译文 $b$ 这是译文~$e'$ 这是译文~$\veca'$, 这是译文 $b \in
  \T$ 这是译文 $\vecv \in \Lambda/\Z^{n}$ 这是译文 (这是译文) 这是译文 $\vecv$ 这是译文  $\veca'$ 这是译文 $\matG^{T} \veca' = \vecv \bmod \Z^{n}$.  这是译文 $\bar{\veca} \in \T_q^n$ 这是译文 $\bar{\vecf} \in \Lambda - \bar{\veca}$, 这是译文~\ref{lem:smoothing}  (这是译文 $r \ge \smootheps(\Lambda)$ 这是译文~\ref{lem:boundonsmoothing}) 这是译文  \begin{align}\label{eq:roundingonelatticeanother}
	  \Pr[\veca = \bar{\veca} \wedge \vecf = \bar{\vecf}] &= q^{-n} \cdot \rho_r(\bar{\vecf}) / \rho_r(\Lambda-\bar{\veca}) \nonumber \\
		& \in C \bracks[\big]{1,\tfrac{1+\epsilon}{1-\epsilon}}\cdot \rho_r(\bar \vecf).
	\end{align}  这是译文 $C=q^{-n}/\rho_{r}(\Lambda)$ 这是译文 $\bar{\veca}$ 这是译文 $\bar{\vecf}$.  这是译文 $\bar \veca, \bar \vecf$ 这是译文 $\bar \veca + \bar \vecf = \bar \vecv$, 这是译文 $\bar{\vecv} \in \Lambda/\Z^{n}$,  \begin{align*}
    \Pr[\vecv = \bar{\vecv}] 
    &\in C \bracks[\big]{1,\tfrac{1+\epsilon}{1-\epsilon}} \cdot
    \rho_{r}(q^{-1} \Z^{n} + \bar{\vecv}).
  \end{align*}  这是译文 $r
  \geq \smootheps(q^{-1} \Z^{n})$ (这是译文~\ref{lem:boundonsmoothing}), 这是译文~\ref{lem:smoothing}  这是译文~$\Pr[\vecv = \bar{\vecv}] \in \bracks[\big]{\tfrac{1-\epsilon}{1+\epsilon},\tfrac{1+\epsilon}{1-\epsilon}} C'$ 这是译文~$C'$ 这是译文~$\bar{\vec{v}}$. 这是译文~\ref{clm:separationdist}, 这是译文~$\veca'$ 这是译文~$1-((1-\eps)/(1+\eps))^2 \le 4\epsilon$ 这是译文.

  这是译文 $A_{q,\vecs,D_{\alpha}}$  这是译文 $A_{q', \matG \vecs, D_{\beta}}$.  这是译文 $(\veca,b=\inner{\veca,\vecs}+e)$, 这是译文 $e
  \gets D_{\alpha}$.  这是译文 $\veca'$ 这是译文 (这是译文)  这是译文~$\T_{q'}^{n'}$.  这是译文 $\overline{\veca'} \in \T_{q'}^{n'}$ 这是译文 $\veca'$,  这是译文 $\bar \vecv = \matG^{T} \overline{\veca'} \bmod \Z^{n}$.  这是译文  \[ b' = \inner{\veca,\vecs} + e + e' = \inner{\overline{\veca'}, \matG \vecs} +
  e + \inner{-\vecf, \vecs} + e' \bmod 1. \]  这是译文~\ref{clm:separationdist} 这是译文~\eqref{eq:roundingonelatticeanother} (这是译文 $\vecf=\bar \vecf$ 这是译文 $\veca = \bar \vecv - \bar \vecf \bmod \Z^n$),  这是译文 $-\vecf$ 这是译文~$1-(1-\eps)/(1+\eps) \le 2 \eps$ 这是译文  $D_{q^{-1}\Z^{n}-\bar \vecv, r}$.  这是译文~\ref{lem:smoothip} (这是译文 $r
  \geq \sqrt{2} \smootheps(q^{-1} \Z^{n})$ 这是译文~$\length{\vecs} \le B$), 这是译文 $\inner{-\vecf,\vecs} + e'$ 这是译文~$6 \eps$  这是译文~$D_{t}$, 这是译文 $t^{2} = r^{2}(\length{\vecs}^{2}+B^{2})$.  这是译文 $e+\inner{-\vecf,\vecs} + e'$ 这是译文~$6 \eps$  这是译文 $D_{(t^2+\alpha^2)^{1/2}}$, 这是译文.
\end{proof}

\begin{comment}

\subsection{Application to Ring-LWE}

\cnote{I think the theorem in this section is a special case of the
  ``ring-LWE embedding'' theorem from GHPS (SCN'12), which says that
  we can embed LWE over a cyclotomic ring with index $m$ into LWE over
  a cyclotomic of index evenly divisible by $m$.  This theorem is the
  special case where $m=1$ (so the cyclotomic ring is just $\Z$).}

\begin{definition}
\label{def:rlwe}
Let~$R_{n,q}$ denote the ring~$\Z_q[x]/(x^n+1)$, with~$n$ a power
of~$2$ and~$q$ an integer. Let~$\chi$ be a distribution
over~$\mathbb{T}_n = \mathbb{T}[x]/(x^n+1)$, and~$\mathcal{D}$ be a
distribution from~$R_{n,q}$. For~$s \in R_{n,q}$, we define the
distribution~$A_{n,q,\chi,s}^R$ as follows: sample~$a \hookleftarrow
U(R_{n,q})$ and~$e \hookleftarrow \chi$; then return~$(a,
\frac{1}{q}(a\cdot s)+e) \in R_{n,q} \times \mathbb{T}_n$.


The decision \emph{Ring Learning with Errors} problem
$\rlwe_{n,q,\chi}(\mathcal{D})$ is as follows: Let~$\vec{s}$ be
sampled from~$\mathcal{D}$; the goal is to distinguish between
arbitrarily many independent samples from~$A_{n,q,\vec{s},\chi}^R$ and
the same number of independent samples from the uniform distribution
over the support of~$A^R_{n,q,\vec{s},\chi}$.
\end{definition}

The above RLWE problem is a special case of the ring learning with
errors problem that was introduced
in~\cite{DBLP:conf/eurocrypt/LyubashevskyPR10}. Since then, it has
proved useful for designing cryptographic schemes with excellent
asymptotic performance. However, all known hardness proofs for \rlwe
rely on the assumed hardness of standard worst-case lattice problems
restricted to the family of so-called ideal lattices. In our present
setup, an ideal lattice is an ideal of~$\Z[x]/(x^n+1)$.
Theorem~\ref{th:reverse} leads to a reduction from \lwe to \rlwe, and
thus provides a hardness proof for \rlwe under standard worst-case
lattice problems over all lattices.

\begin{theorem}
\label{th:rlwe}
For any~$q, \chi, \mathcal{D}$, there exists a polynomial-time
reduction from~$\lwe_{1,q,\chi}(\mathcal{D})$
to $\rlwe_{n,q,\chi'_n}(\mathcal{D}_n)$ where~$\chi'_n = \chi'+\chi' x
+ \cdots + \chi' x^{n-1}$, $\chi'$ is the sum of $n$ independent
copies of~$\chi$ and~$\mathcal{D}_n = \mathcal{D} + \mathcal{D} x +
\cdots + \mathcal{D} x^{n-1}$.
\end{theorem}

\begin{proof}
  For~$s_0,\ldots,s_{n-1} \in \Z_p$, we define the
  distribution~$A_{n,q,\chi,(s_i)_i}$: Sample~$a \hookleftarrow
  U(\Z_q)$ and~$e_0,\ldots,e_{n-1} \hookleftarrow \chi$, and
  return~$(a, \frac{1}{p}(a\cdot
  s_0)+e_0,\ldots,\frac{1}{p}(as_{n-1})+e_{n-1}) \in \Z_p \times
  \mathbb{T}^n$.  Using a standard hybrid argument, we see
  that~$\lwe_{1,q,\chi}(\mathcal{D})$ reduces to distinguishing between
  arbitrarily many independent samples from~$A_{n,q,\chi,(s_i)_i}$ and
  the same number of independent samples from the uniform distribution
  over the support of~$A_{n,q,\chi,(s_i)_i}$. We describe an efficient
  mapping of~$n$ elements~$(a_i, (b_{ij})_{0\leq j<n})_{0 \leq i<n}$
  from~$\Z_p \times \mathbb{T}^n$ to a single element in~$R_q \times
  \mathbb{T}_n$ such that the $A_{n,q,\chi,(s_i)_i}$ distribution
  (resp.\ uniform distribution over the support
  of~$A_{n,q,\chi,(s_i)_i}$) is mapped to the~$A_{n,q,\chi_n,s}^R$
  distribution (resp.\ uniform distribution over the support
  of~$A_{n,q,\chi_n,s}^R$), where~$s = s_0+s_1x + \cdots
  +s_{n-1}x^{n-1}$. The mapping is as follows:
\begin{itemize}
\item Define~$a = a_0 + a_1x + \cdots + a_{n-1}x^{n-1} \in R_{n,q}$.
\item Define~$b = \sum_{0\leq k<2n-1} (\sum_{i+j = k} b_{ij}) x^k \bmod (x^n+1)$.
\item Return~$(a,b)$.  
\end{itemize}
Clearly, the reduction runs in polynomial-time, and the uniform
distribution is mapped to the uniform distribution. Finally, in case
the input~$b_{ij}$'s are of the form~$\frac{1}{p} (a_i\cdot s_j) +e_j$
for all~$i,j$, then all the coefficients of~$b - \frac{1}{p}(a\cdot s)
\in \mathbb{T}_n$ are sums of~$n$ independent samples from~$\chi$.
\end{proof}

As a corollary of Theorems~\ref{th:reverse} and~\ref{th:rlwe}, there
exists a (classical) polynomial-time reduction from standard
worst-case problems over lattices with polynomial approximation
factors to~$\rlwe_{n,p^n,D_{\sqrt{n}\alpha}}(U(R_{n,p^n}))$.

\end{comment}

\begin{comment}
\begin{verbatim}
Also, we were wondering about a possible relationship with
   http://www.di.ens.fr/~lyubash/papers/subsetsumcrypto.pdf
(digit-wise expansion of a large integer, along columns rather than rows).

Notation: \vec* denotes an n-dimensional vector mod p (i.e., in Zp^n),
while capital letters (A,S) denote scalar values mod p^n (i.e., in
Z_{p^n}).

The basic approach is to map LWE samples (\vecr,b) in Zp^n \times Zp
to (A, b*p^{n-1} + noise) in Z_{p^n} \times Z_{p^n}, where A is chosen
from r using some *randomized* encoding (not deterministically).  Then
the amount by which b*p^{n-1} differs from the value A*S will be a
centered Gaussian, and we only need to add a small amount of extra
noise (or maybe none at all?) to smooth it out.

Details:

For A,S in Z_p^n we can write A = sum_i p^i * Ai for *integers* ai
(not necessarily in {0,...,p-1}), and similarly for S.

Then the product A*S mod p^n can be written as \veca^t \matP \vecs mod
p^n, where

\veca = (a_{n-1}, ..., a_0),   \vecs = (s0, s1, ..., s_{n-1}),  and
\matP is the fixed lower-triangular Toeplitz matrix with
\vecp=(p^{n-1}, p^{n-2}, ..., p^0) down its leftmost column.

Ideally, we would then like to choose \veca (defining A) to satisfy

  \veca^t \matP = p^{n-1} * \vecr^t    (mod p^n)

where \vecr is the "left-hand side" of the LWE challenge.

If we could do this, then the product A*S would be exactly p^{n-1} *
<r,s>.  We are given the noisy value b=<r,s>+error, so we could just
plug in p^{n-1}*b (mod p^n) to complete the reduction.  Notice that
this would keep the error rate exactly the same, and moreover the
error term would be divisible by p^{n-1}, which seems like too much to
ask.... and it is.

The problem, of course, is that it is usually not possible to solve
the above system, because \matP is highly non-invertible mod p^n.  (It
doesn't even help that the RHS is a multiple of p^{n-1}.)

Luckily, it's enough to approximately solve the system, as long as the
"error" \veca^t \matP - p^{n-1} * \vecr^t is a centered Gaussian and
\vecs is short.  And the structure of \matP makes this easy to do --
Oded's formulas are the answer.

[Chris note: the formulas come from an earlier email, with slightly
different notation.  But they are easy to work out: the entries of
\veca are chosen iteratively from discrete Gaussians over the
integers, with appropriate centers determined by the previously chosen
entries of \veca.]

Another way of going about it is to multiply both sides by \matP^-1
**over the integers**, which has a very simple form (exercise for
reader).  Then we end up with a fractional RHS which we can just
randomized-round to get an integral \veca.  Doing it this way causes
the "error" \veca^t \matP - p^{n-1} * \vecr^t to be a non-spherical
but "nearly spherical" Gaussian (because \matP is almost orthogonal),
which would suffice for our purposes.  It's not clear that this
non-adaptive solution has any important advantage here, though.

(Note: these two approaches correspond to randomized versions of
Babai's nearest-plane and "simple rounding" algorithms, respectively.)

\end{verbatim}
\end{comment}



%%% Local Variables: 
%%% mode: latex
%%% TeX-master: "lwehard"
%%% End: 
